\section{Details}\label{app:details}
% Material moved from the main text (final review, 30 Sep 2026). Sources: claims.md v3.9
% (C3, C5, C6, B3, A7, S1-S6), docs/PREREGISTRAZIONE-scopesweep.md (derivation of the prior
% prediction), tools/fig02b_data.py (rates of the exhaustive series),
% tools/honeypot_aggregate.py (signatures: presence of the headers only).

\subsection{Honeypot}\label{app:honeypot}
The public measurement site is \texttt{theslowshelf.org}.
The data used in this paper cover the window [12 August, 22 September 2026) UTC; traffic
after the publication of this paper, which names the site, is not comparable with it.
Its web server logs, for each request, the time, URL,
status and user agent, an identifier of the TCP connection with the request's position in
it, and the Web Bot Auth headers \texttt{Signature-Agent} and
\texttt{Signature-Input}.\claim{C6}
Clients are classed by their declared user agent into six classes; the \emph{agent} class
matches fetchers that act for a user, such as \texttt{ChatGPT-User}, \texttt{Claude-User}
and Perplexity.
The declaration is not verified.
All aggregates come from one fingerprinted snapshot of the logs, with 520,871 requests in
that window; the handling of client
addresses is described in \cref{app:ethics}.\claim{C6}

The windowed working set is computed for the clients with at least 5 requests.\claim{C3}
Its 143\,s window is the characteristic time of the testbed cache at the smallest agentic
setting of the earlier runs: the capacity, then estimated at 5,274 objects, divided by the
origin miss rate of about 37 req/s at scope 0.02; with the measured capacity of 5,263
objects it would be 142\,s, a difference of 0.2\% that we do not recompute.\claim{C3}
The distribution of the windowed working set across clients, and the counts behind the
contiguity comparison, are given with \cref{fig:a-honeypot}.\claim{C3}\claim{C5}
Of the requests with signature headers, 19,934 of 520,871, all but two come from generic
bots.\claim{C6}

\subsection{The prior characteristic-time prediction}\label{app:tc}
Under the characteristic-time approximation~\cite{fagin1977,che2002,fricker2012}, an object
in an LRU cache of capacity $C$ survives if it is requested again within
$T_C = C/M$, where $M$ is the rate of misses.
If a class sends $\lambda$ requests per second spread uniformly over $W$ objects, each
object is requested at rate $\lambda/W$ and the class's miss ratio is
$\exp(-\lambda T_C / W)$.
The prior prediction used this form with $T_C = 143$\,s and the reachable sets as then
computed, 3,052 and 10,172 objects, so that the elasticity
$\mathrm{miss}(12)/\mathrm{miss}(36) = \exp(24\,T_C/W)$ gives the two predicted values of
\cref{sec:m-tc}.\claim{B3}
\paragraph{The post-hoc calculation.}
% claims.md B7; docs/RISULTATO-che-posthoc.md; tools/che_posthoc.py.
The calculation of \cref{tab:tc} keeps the approximation and replaces its inputs with those
of our design.\claim{B7}
Each chapter $j$ receives a request rate $\mu_j$ from the human class, whose rank
$r = \lfloor N^u \rfloor$ has probability $\log((k+1)/k)/\log N$, and from the agentic class,
whose session starts have probability $((k+1)/S)^{0.4} - (k/S)^{0.4}$ over the first
$S = \lfloor N \cdot \mathrm{scope} \rfloor$ ranks, each session adding a third of its rate to
three consecutive chapters.\claim{B7}
The exhaustive class sends $\alpha\lambda$ requests per second over the $N$ chapters.
In version (A) it is Poisson traffic like the others, so a chapter is a hit with probability
$1 - e^{-(\mu_j + \alpha\lambda/N) T}$.\claim{B7}
In version (B) it returns to each chapter every $P = N/(\alpha\lambda)$ seconds, as in the
corrected generator; with $T < P$, the chapter is in the cache with probability
$1 - e^{-\mu_j T}(1 - T/P)$, which is also the hit probability of the other classes, and an
exhaustive request hits with probability $1 - e^{-\mu_j T}$.\claim{B7}
$T$ solves the capacity equation $\sum_j s_j\, o_j(T) = C$, where $o_j$ is the occupancy
probability of chapter $j$, $s_j$ its size from the reconstructed responses plus a fixed
512-byte overhead, and $C$ the 128\,MiB of the cache.\claim{B7}
The origin rate is the sum of the miss rates, and the marginal is computed as for the
measurements.
No parameter is fitted to the measurements: the sizes and the overhead are those fixed in
advance for the simulator (\cref{app:sim}), and $T$ follows from the capacity.\claim{B7}

\subsection{Simulator}\label{app:sim}
The first phase checked the simulator against the testbed, with 20 seeds per point and
tolerances fixed in advance.\claim{S1}
It reproduces the per-class miss ratios (18 of 18 within 0.02 of the testbed) and the origin
rate (within 3\%, from $+0.37\%$ to $+0.92\%$).\claim{S1}
The simulated agentic marginal at scope 0.02 is $-0.0451 \pm 0.0033$, against
$-0.0405 \pm 0.0061$ measured, and the simulated overlap effect is $+0.882$ and $+0.839$
origin req/s, against $+0.882$ and $+0.872$; the comparison used the shared traversal
counter on both sides.\claim{S1}
The simulated origin rate is slightly higher than the measured one at all five points
($+0.4\%$ to $+0.9\%$), and the simulated miss ratios almost everywhere, within the
tolerance.\claim{S1}
Object sizes are the reconstructed chapter responses plus a fixed 512-byte overhead, fixed
in the pre-registration of the first phase.\claim{S1}\claim{C8}
\Cref{tab:sim} gives the outcome of the second phase, with the corrected traversal, for all
eight combinations.\claim{S2}\claim{S3}\claim{S4}\claim{S5}
Every value in it is simulated; policies other than LRU were not validated on the testbed.

\begin{table}[h]
\caption{Second phase of the simulator (simulated, not measured): four eviction policies
and two session models (the generator's, and sessions calibrated on the honeypot). Marginal
at scope 0.02 and net effect of adding the agentic class from 0 to 36 req/s, in origin
requests per request and per second.}
\label{tab:sim}
\centering\scriptsize
\begin{tabularx}{\textwidth}{@{}ll>{\raggedleft\arraybackslash}X>{\centering\arraybackslash}X>{\centering\arraybackslash}X>{\centering\arraybackslash}X>{\raggedleft\arraybackslash}X@{}}
\toprule
Policy & Sessions & Marginal at 0.02 & Depends on reachable set & Sign change & Overlap effect at 36 req/s & Net 0 $\to$ 36 \\
\midrule
LRU       & generator  & $-0.031$ & yes & yes & yes & $+1.20$ \\
LRU       & calibrated & $-0.017$ & yes & yes & yes & $+1.43$ \\
SIEVE     & generator  & $-0.027$ & yes & yes & no  & $-0.42$ \\
SIEVE     & calibrated & $-0.025$ & yes & yes & no  & $-0.42$ \\
S3-FIFO   & generator  & $-0.046$ & yes & yes & no  & $-1.05$ \\
S3-FIFO   & calibrated & $-0.043$ & yes & yes & no  & $-0.99$ \\
W-TinyLFU & generator  & $+0.004$ & no  & no  & no  & $+0.56$ \\
W-TinyLFU & calibrated & $+0.001$ & yes & no  & no  & $+0.58$ \\
\bottomrule
\end{tabularx}
\end{table}

\subsection{The secondary exhaustive series}\label{app:a7}
To check that the exhaustive class is insensitive to its own volume, and not only to the
agentic one, a campaign of 16 September 2026 varied the exhaustive volume from 0 to 42 req/s
with the human and agentic classes fixed at 55 and 12 req/s.\claim{A7}
It used the shared mapping, a measurement window of 1,211\,s and 3 repetitions per point,
and the shared traversal counter.\claim{A7}
The configured exhaustive rates are 0, 13.9968, 27.9965 and 41.9977 req/s; their provenance
is stated in \cref{sec:limitations}.\claim{A7}
The marginal cost of the exhaustive class is $0.961 \pm 0.009$, $0.997 \pm 0.010$ and
$0.979 \pm 0.004$ over the three steps, and its miss ratio rises from 0.791 to 0.861
between 14 and 42 req/s.\claim{A7}
